% GENERATED FILE. Edit canonical lessons, then run npm run generate:books.
% canonical-source-hash: fnv1a64:8096bf9cff244712
% canonical-lessons: PT-C261-acho-que, PT-C261-porque, PT-C261-por-isso, PT-C261-revisao-opiniao

\chapter{Acho Que --- What I Think, and Why}
\label{ch:pt-what-i-think-and-why}
\hlchaptermodality{261}

\begin{chapteropening}
\textbf{By the end of this chapter:} \emph{I can give my opinion with acho que or para mim, give a reason with porque, and run it forward with por isso or então.}

Complete PT-C261-revisao-opiniao: an opinion, its reason and its consequence, said in either order.
\end{chapteropening}

\section[acho que, para mim]{acho que — I find that, which is to say, I think}

\label{lesson:PT-C261-acho-que}



\begin{quote}
Say \textquotedblleft{}I think so\textquotedblright{} with the verb chapter nineteen gave you. (\emph{Penso que sim.})

Giving reasons starts with what you think.
\end{quote}

\subsection*{You'll want to know: acho que}
\textbf{acho que} — \textquotedblleft{}I think that\textquotedblright{}. It is the everyday opinion opener, in Portugal as much as in Brazil, from \textbf{achar}, a regular \emph{-ar} verb meaning \textquotedblleft{}to find\textquotedblright{}.

\begin{quote}
\textbf{Acho que} o museu é interessante. — I think the museum is interesting. \textbf{Acho que} o comboio é caro. — I think the train is expensive. \textbf{Acho que sim.} — I think so. · \textbf{Acho que não.} — I don't think so.
\end{quote}

To hand the opinion back, use the same verb: to a friend, \textbf{E tu, o que achas?} — \textquotedblleft{}and you, what do you think?\textquotedblright{} — and to \emph{o senhor} or \emph{a senhora}, \textbf{O que acha?}

\emph{Penso que} says the same thing, a little more carefully. After either one the verb stays plain: \emph{acho que é verdade}. Only a denial, \emph{não acho que…}, makes careful Portuguese bend the verb, which is a later chapter's work.

\begin{grammarlens}[title={Grammar Lens: para mim}]
\textbf{para mim} — \textquotedblleft{}for me\textquotedblright{}, and so \textquotedblleft{}to my mind\textquotedblright{}:

\begin{quote}
\textbf{Para mim}, o comboio é caro, mas o museu é interessante. — To my mind, the train is expensive, but the museum is interesting.
\end{quote}

\emph{Mim} is the form \emph{eu} takes after \emph{para}, and \emph{tu} becomes \emph{ti}, so \textbf{para ti?}, \textquotedblleft{}what about you?\textquotedblright{}, hands the view back in two words. \emph{Para mim} puts your own standpoint first; \emph{acho que} puts the opinion first.
\end{grammarlens}

\begin{cousinweb}
\textbf{achar} is inherited from Latin \emph{afflāre}, \textquotedblleft{}to breathe upon\textquotedblright{} (EtymDB 2.1), a hunter's word for catching a scent. To find a thing was to pick up its scent, so \emph{acho que} is, at root, \textquotedblleft{}I find that\textquotedblright{}, as in English \textquotedblleft{}I find it expensive\textquotedblright{}.
\end{cousinweb}

\subsection*{Your turn}
Say these aloud:

\begin{itemize}
\raggedright
  \item \emph{Acho que o museu é interessante.}
  \item hand it back — \emph{E tu, o que achas?}
  \item both answers — \emph{Acho que sim.} then \emph{Acho que não.}
  \item \emph{Para mim, o comboio é caro.}
\end{itemize}

\begin{tcolorbox}[breakable,colback=teal!4,colframe=teal!35!black,title={Before you move on}]
What does \emph{achar} mean, plainly? (\textbf{To find.}) Say \textquotedblleft{}I don't think so\textquotedblright{}. (\emph{\textbf{Acho que não.}}) How do you ask a friend what they think? (\emph{\textbf{O que achas?}}) And what does \emph{eu} become after \emph{para}? (\emph{\textbf{Mim}} — \emph{para mim}.)
\end{tcolorbox}
\section[porque]{porque — because, at last}

\label{lesson:PT-C261-porque}



\begin{quote}
Say \textquotedblleft{}I think so\textquotedblright{}. (\emph{Acho que sim.}) Then join two facts with the joiner from the writing chapters: \textquotedblleft{}it was cold, that's why I didn't go out\textquotedblright{}. (\emph{Fazia frio, por isso não saí.})

When you wrote four lines, this book had no word for \emph{because}. Here it is.
\end{quote}

\subsection*{You'll want to know: porque}
\textbf{porque} — \textquotedblleft{}because\textquotedblright{}. It goes in front of the reason:

\begin{quote}
\emph{Não saí \textbf{porque} chovia.} — I didn't go out because it was raining. \emph{Apanhei o autocarro \textbf{porque} fazia frio.} — I took the bus because it was cold.
\end{quote}

Notice the tenses keep their jobs. \emph{Não saí} and \emph{apanhei} are events, in the perfeito; the reasons, \emph{chovia} and \emph{fazia frio}, are scene, in the imperfeito. A reason in a story is very often part of the scene.

\begin{grammarlens}[title={Grammar Lens: porque é que…?}]
The same word asks the question, and European Portuguese likes to ask it with \textbf{é que} after it, \textquotedblleft{}is it that\textquotedblright{}:

\begin{quote}
— \emph{\textbf{Porque é que}} estudas português? — \emph{\textbf{Porque}} gosto de Portugal. — Why are you learning Portuguese? — Because I like Portugal.
\end{quote}

When the word stands alone, or ends the sentence, it takes an accent and the stress moves: \emph{Não vou.} — \emph{\textbf{Porquê?}} — \textquotedblleft{}I'm not going.\textquotedblright{} — \textquotedblleft{}Why?\textquotedblright{}

Brazil writes the question as two words, \emph{por que}, and the lone one as \emph{por quê}; the answer is \emph{porque} everywhere.
\end{grammarlens}

\begin{cousinweb}
\textbf{porque} is two words grown together: \textbf{por}, \textquotedblleft{}for\textquotedblright{}, and \textbf{que}, \textquotedblleft{}what, that\textquotedblright{} (EtymDB 2.1 lists the compound, with Italian \emph{perché} beside it). \textbf{que} is inherited from Latin \emph{quid}, \textquotedblleft{}what\textquotedblright{} (EtymDB 2.1), and \textbf{por} is usually traced to Latin \emph{prō}, \textquotedblleft{}for, on behalf of\textquotedblright{}. So the question and the answer are the same two words: \textquotedblleft{}for what\textquotedblright{}, and \textquotedblleft{}for this\textquotedblright{}.
\end{cousinweb}

\subsection*{Your turn}
Say these aloud:

\begin{itemize}
\raggedright
  \item \emph{Não saí porque chovia.}
  \item \emph{Apanhei o autocarro porque fazia frio.}
  \item the question — \emph{Porque é que estudas português?}
  \item your own answer — \emph{Porque …}
\end{itemize}

\begin{tcolorbox}[breakable,colback=teal!4,colframe=teal!35!black,title={Before you move on}]
Where does \emph{porque} go: before the reason, or before the result? (\textbf{Before the reason.}) Ask \textquotedblleft{}why are you learning Portuguese\textquotedblright{} the European way. (\emph{\textbf{Porque é que estudas português?}}) When does it take an accent? (When it stands \textbf{alone} or ends the sentence: \emph{\textbf{porquê}}.)
\end{tcolorbox}
\section[por isso, então]{por isso, então — the arrow that points forward}

\label{lesson:PT-C261-por-isso}



\begin{quote}
Give an opinion with a reason: \textquotedblleft{}I think it's going to rain, because there's no sun.\textquotedblright{} (\emph{Acho que vai chover, porque não há sol.})

\emph{Porque} runs from the claim back to its reason. This lesson runs the other way.
\end{quote}

\begin{grammarlens}[title={Grammar Lens: the two arrows}]
The same two facts, in opposite orders:

\begin{quote}
Apanhei o autocarro \textbf{porque} chovia. — I took the bus \textbf{because} it was raining. Chovia, \textbf{por isso} apanhei o autocarro. — It was raining, \textbf{so} I took the bus.
\end{quote}

\textbf{Porque} points \textbf{back}, from what happened to its cause. \textbf{Por isso} points \textbf{forward}, from the cause to what followed; its \emph{isso}, \textquotedblleft{}that\textquotedblright{}, reaches back to the cause, which is why the cause has to come first. Which one you reach for depends only on which fact you want to say first, and that is a choice about emphasis, not about grammar.

The tenses keep their jobs. \emph{Chovia} is the scene, the reason; \emph{apanhei} is the event, the consequence.
\end{grammarlens}

\begin{grammarlens}[title={Grammar Lens: então and portanto}]
\textbf{então} can stand in the same place: \emph{Chovia, \textbf{então} apanhei o autocarro.} It is lighter and more conversational, the \textquotedblleft{}so\textquotedblright{} of a friend telling you about their day. \textbf{portanto}, which you also have, is the formal one: a letter, a report, an argument.

\begin{quote}
spoken: \emph{então} · neutral: \emph{por isso} · written: \emph{portanto}
\end{quote}

All three point forward. \emph{Então} also keeps its first meaning, \textquotedblleft{}then, at that time\textquotedblright{}, so in a story it can do both jobs at once.
\end{grammarlens}

\begin{cousinweb}
\textbf{então} is inherited from Latin \emph{in tunc}, \textquotedblleft{}in that moment\textquotedblright{} (EtymDB 2.1; Lewis \& Short, under \emph{tunc}): a word about time that became a word about consequence, exactly as English \textquotedblleft{}then\textquotedblright{} did. \textbf{isso} is inherited from Latin \emph{ipsum}, \textquotedblleft{}itself, that very thing\textquotedblright{} (EtymDB 2.1), so \emph{por isso} is \textquotedblleft{}for that very thing\textquotedblright{}.
\end{cousinweb}

\subsection*{Your turn}
Say these aloud:

\begin{itemize}
\raggedright
  \item back — \emph{Apanhei o autocarro porque chovia.}
  \item forward — \emph{Chovia, por isso apanhei o autocarro.}
  \item the same, to a friend — \emph{Chovia, então apanhei o autocarro.}
  \item \emph{Para mim, o comboio é caro, por isso apanho o autocarro.}
\end{itemize}

\begin{tcolorbox}[breakable,colback=teal!4,colframe=teal!35!black,title={Before you move on}]
Which way does \emph{porque} point, and which way \emph{por isso}? (\emph{Porque} \textbf{back} to the cause; \emph{por isso} \textbf{forward} to the consequence.) Which of the three forward words is the spoken one, and which the written one? (\emph{\textbf{Então}} spoken; \emph{\textbf{portanto}} written.) What did Latin \emph{in tunc} mean? (\textbf{At that moment.})
\end{tcolorbox}
\section[(review)]{Review — what you think, why, and what follows}

\label{lesson:PT-C261-revisao-opiniao}



\begin{quote}
Someone asks \emph{O museu é interessante?} Answer \textquotedblleft{}I think so\textquotedblright{}. (\emph{Acho que sim.}) Then ask them back. (\emph{E tu, o que achas?})
\end{quote}

\begin{grammarlens}[title={What you've built}]
Three moves, and you own words for each. An \textbf{opinion}: \emph{acho que}, or \emph{para mim}. A \textbf{reason}, pointing back: \emph{porque}. A \textbf{consequence}, pointing forward: \emph{por isso}, or \emph{então}.

\begin{quote}
\emph{Acho que o museu é interessante, porque é muito antigo.} \emph{O comboio é caro, por isso apanho o autocarro.}
\end{quote}

The first sentence runs claim, then cause. The second runs cause, then consequence. Put them side by side and you can argue a small point in either direction, which is most of what giving reasons is.
\end{grammarlens}

\subsection*{Your turn}
Say these aloud:

\begin{itemize}
\raggedright
  \item \emph{Acho que o museu é interessante, porque é muito antigo.}
  \item the same facts, forward — \emph{O museu é muito antigo, por isso é interessante.}
  \item \emph{O comboio é caro, então apanho o autocarro.}
  \item a reason in the past — \emph{Chovia, por isso não saí.}
  \item answer \emph{Porque é que estudas português?} with a \emph{porque} and a \emph{mas}
\end{itemize}

\begin{tcolorbox}[breakable,colback=teal!4,colframe=teal!35!black,title={Before you move on}]
Name the three moves and a word for each. (Opinion \emph{\textbf{acho que}}, reason \emph{\textbf{porque}}, consequence \emph{\textbf{por isso}} or \emph{\textbf{então}}.) Say \textquotedblleft{}to my mind\textquotedblright{}. (\emph{\textbf{Para mim.}}) In \emph{chovia, por isso não saí}, which half is the scene? (\emph{\textbf{Chovia}} — the reason.)
\end{tcolorbox}
